%%%%%%%%%%%%%%%%%%%%%%%%%%%%%%%%%%%%%%%%%%%%%%%%%
\section{幂次计数}
\label{sec:powercounting}

\subsection{表面紫外发散}

为了辨认任何对准 PDF 有贡献的费曼图的全部紫外发散，我们从一组一般图形出发：向低阶费曼图添加一条额外的胶子来构造新图形，然后考察这条附加胶子会如何改变原图形的紫外发散。具体地，我们引入并研究发散指标的变化 $\Delta\omega_3$ 与 $\Delta\omega_4$，分别对应相应圈的三维积分与四维积分。准 PDF 可重正的一个充分条件是所有情形都满足 $\Delta\omega_3\le 0$ 且 $\Delta\omega_4\le 0$，但正如我们将要论证的，这不是必要条件。我们把讨论分为五种不同的情形：

%%%%%%%%%%%%%%%%%%%%%%%%%%%%%%%
\begin{figure}[t]
\begin{center}
\includegraphics[width=3in]{addonegluon}
\caption{\label{fig:addonegluon}
一些圈图及可能的附加胶子（以虚线表示的连接），它们可以在更高圈阶对准 PDF 有贡献。}
\end{center}
\end{figure}
%%%%%%%%%%%%%%%%%%%%%%%%%%%%%%%

\textbf{情形 I：}胶子只有一端连接到一个圈的部分子线上，如图~\ref{fig:addonegluon}(a)。此时圈内多了一条传播子和一个 QCD 顶点，导致 $\Delta\omega_3=-1$ 与 $\Delta\omega_4=-1$。于是线性发散可以变为对数发散，而对数发散的图则变为有限。

\textbf{情形 II：}胶子只有一端连接到一个圈的规范链接上，如图~\ref{fig:addonegluon}(b)。此时多了一个坐标空间中 $z$ 方向的积分，它不改变三维积分的发散度，但使四维积分的发散度减一，即 $\Delta\omega_3=0$ 且 $\Delta\omega_4=-1$。

\textbf{情形 III：}胶子的两端都连接到一个圈的部分子线上，如图~\ref{fig:addonegluon}(c)。此时多了三条传播子、两个 QCD 顶点以及圈的四个额外动量积分，导致 $\Delta\omega_3=-1$ 与 $\Delta\omega_4=0$。

\textbf{情形 IV：}胶子的一端连接到一个圈的部分子线上，另一端连接到该圈的规范链接上，如图~\ref{fig:addonegluon}(d)。此时多了两条传播子、一个 QCD 顶点、一个坐标空间中 $z$ 方向的积分以及圈的四个额外动量积分，导致 $\Delta\omega_3=0$ 与 $\Delta\omega_4=0$。

\textbf{情形 V：}胶子的两端都连接到一个圈的规范链接上，如图~\ref{fig:addonegluon}(e,f)。此时多了一条传播子、两个坐标空间中 $z$ 方向的积分以及圈的四个额外动量积分。结果是 $\Delta\omega_3= 1$ 与 $\Delta\omega_4=0$。

对上述幂次计数规则作几点说明。
1) 我们只关注整体发散而不关注子发散，因为子发散可以在重正化步骤中以森林减除处理。2) 这里讨论的发散指标变化只针对表面发散：这意味着即使幂次计数表明某图发散，由于其他方面的考虑它也可能实际不发散；但反过来，若幂次计数表明某图有限，则它必定是紫外有限的。

\textbf{情形 V} 对重正化是危险的，因为 $\Delta\omega_3>0$ 意味着潜在紫外发散拓扑的数目不再有限，从而有限个重正化常数可能不足以消除全部紫外发散。然而如前所指，上述幂次计数规则只针对表面发散。下一小节将证明：只要对其任一圈动量只作三维积分，准 PDF 的所有费曼图都没有整体紫外发散。因此，只有 $\Delta\omega_4$ 对辨认真实的紫外发散拓扑有意义。

%%%%%%%%%%%%%%%%%%%%%%%%%%%%%%%
\subsection{三维积分的有限性}

为了证明对准 PDF 有贡献的费曼图的三维积分不会产生真实的紫外发散，我们考虑``规范链接不可约''（GLI）图。类似于单粒子不可约（1PI）图的概念，GLI 图是指把所有规范链接剪断（或移除）之后仍然连通的图。例如，图~\ref{fig:oneloop} 中的图 (a)、(b)、(c) 不是 GLI 图，而图 (d) 是；类似地，图~\ref{fig:addonegluon} 中的图 (a)、(c)、(d) 是 GLI 图，而图 (b)、(e)、(f) 不是。研究准 PDF 的重正化时，由于 QCD 拉氏量的重正化众所周知，只需研究图~\ref{fig:connected} 所示的一般 GLI 图的紫外性质即可，其中连接到规范链接的虚线既可以是胶子也可以是夸克（如果它位于规范链接的端点上）。

%%%%%%%%%%%%%%%%%%%%%%%%%%%%%%%
\begin{figure}[h]
\begin{center}
\includegraphics[width=1in]{strongconnected}
\caption{\label{fig:connected}
一般的``规范链接不可约''（GLI）图，其中 $l_0=q-l_1-\cdots-l_n$。}
\end{center}
\end{figure}
%%%%%%%%%%%%%%%%%%%%%%%%%%%%%%%
在图~\ref{fig:connected} 中，假设流入圆团的动量为 $q$，$l_1,\cdots,l_n$ 为 $n$ 个圈动量，而 $l_0=q-l_1-\cdots-l_n$ 由动量守恒确定。注意，这种动量指派对 GLI 图总是可行的。于是，图~\ref{fig:connected} 的一般图可以表示为
\begin{align}\label{eq:connected}
e^{i q_{z} r_0} \prod_{j=1}^n \int_{r_{j-1}+a}^{r_{max}-a}d r_j \int \frac{d^4l_j}{(2\pi)^4} e^{i l_{jz} (r_j-r_0)} \mathcal{M}(q,l_1,\cdots,l_n),
\end{align}
其中 $\mathcal{M}(q,l_1,\cdots,l_n)$ 表示圆团连同连接到规范链接的那 $n+1$ 条线的传播子。GLI 图可以由图~\ref{fig:oneloop} 或图~\ref{fig:oneloopg} 中的一圈图结合上一小节 \textbf{情形 I、III} 和 \textbf{IV} 的插入构造而成，因此无论对各圈动量施加三维积分还是四维积分，它们的整体表面紫外发散指标都是 $\omega\leq1$。

现在研究只对 $l_j$ 作三维积分、而对其他圈动量作三维或四维积分的情形。$\mathcal{M}(q,l_1,\cdots,l_n)$ 对 $l_j$ 有两类依赖：一类是分子中的多项式依赖，对此我们可以简单地把 $l_j$ 分解为 $\bar{l}_{j}$ 与 $l_{jz}$；另一类是形如 $1/(l_j+k)^2$ 的传播子依赖，其中 $k$ 可以为零或依赖其他圈动量。可以把该传播子分解为
\begin{align}\label{eq:decomZ}
\begin{split}
\frac{1}{(l_j+k)^2}=& \frac{1}{\Delta-2k_z l_{jz}-l_{jz}^2}\\
=& \frac{1}{\Delta}+ \frac{2k_z l_{jz}}{\Delta^2}+\frac{(\Delta+4k_z^2+2k_zl_{jz}) l_{jz}^2}{(\Delta-2k_z l_{jz}-l_{jz}^2)\Delta^2},
\end{split}
\end{align}
其中 $\Delta=(\bar{l}_j+\bar{k})^2-k_z^2$ 与 $l_{jz}$ 无关。根据量纲计数，分子中每多出一个 $l_{jz}$，就会把三维积分的发散指标压低一个单位。因此，既然我们考虑的是 $l_j$ 三维积分的紫外发散，式~\eqref{eq:decomZ} 的最后一项可以被安全地忽略。其余各项把 $l_{jz}$ 依赖从 $\mathcal{M}(q,l_1,\cdots,l_n)$ 中因子化出来，这些项的潜在紫外发散正比于
\begin{align}
\int d l_{jz} e^{i l_{jz} (r_j-r_0)}  l_z^m \propto \delta^{(m)}(r_j-r_0),
\label{eq:deltam}
\end{align}
其中 $m$ 为非负整数。由于按式~\eqref{eq:connected} 的定义 $r_j$ 不能等于 $r_0$，式~(\ref{eq:deltam}) 中的 $\delta^{(m)}(r_j-r_0)$ 在取 $a\to 0$ 极限之前就消失了。因此，图~\ref{fig:connected} 中先积掉 $\bar{l}_j$ 及其他圈动量、固定 $l_{jz}$ 而得到的紫外发散，最终在对 $l_{jz}$ 积分后消失。

总之，GLI 图的整体紫外发散只来自所有圈动量都很大的区域。此外，正是式~\eqref{eq:decomZ} 的第三项负责 GLI 图四维积分的真实紫外发散。



%%%%%%%%%%%%%%%%%%%%%%%%%%%%%%%
\begin{figure}[h]
\begin{center}
\includegraphics[width=1.3in]{disconnected}
\caption{\label{fig:disconnect}
由两个``规范链接不可约''（GLI）子图构成的图。}
\end{center}
\end{figure}
%%%%%%%%%%%%%%%%%%%%%%%%%%%%%%%

现在研究图~\ref{fig:disconnect} 所示由两个 GLI 子图构成的图。这个非 GLI 图既可以由图~\ref{fig:twoloop} 中的图生成，也可以由上一小节 \textbf{情形 II} 和 \textbf{V} 的插入生成。按照幂次计数规则，此图的整体表面紫外发散指标为 $\omega\leq2$。提取整体紫外发散时，可以对每个 GLI 子图沿用上面的讨论，发现只要固定任一圈动量的 $z$ 分量，合并图的整体紫外发散就会消失。这一结论容易推广到由更多 GLI 子图构成的任何图。


%%%%%%%%%%%%%%%%%%%%%%%%%%%%%%%
\begin{figure}[h]
\begin{center}
\includegraphics[width=1.1in]{twoloopbase}
\caption{\label{fig:twoloop}
基于幂次计数规则以及图~\ref{fig:oneloop} 与图~\ref{fig:oneloopg} 中的基本构件，可能给夸克准 PDF 带来紫外发散贡献的一个两圈图。}
\end{center}
\end{figure}
%%%%%%%%%%%%%%%%%%%%%%%%%%%%%%%

由此我们得出结论：凡对准 PDF 有贡献的图，其整体紫外发散只能来自所有圈动量都很大的区域。这一发现的直接推论是：评估上一小节讨论的五种插入以辨认紫外发散拓扑时，只需考察 $\Delta\omega_4$ 的取值。既然我们发现对所有情形都有 $\Delta\omega_4\le 0$，可能对准 PDF 有贡献的紫外发散费曼图拓扑应当只有有限个，我们将在下一小节给出。

最后说明一下准 PDF 与 PDF（后者具有来自三维积分的紫外发散）之间紫外行为的关键差异。如第 \ref{sec:oneloop} 节末尾所指出的，PDF 的紫外发散来自对圈动量另一组分量的三维积分，比如 $l_-$ 与 $\vec{l}_\perp$。当然，可以对传播子做与式~\eqref{eq:decomZ} 类似的分解：
\begin{align}\label{eq:decomPlus}
\begin{split}
\frac{1}{(l+k)^2}=& \frac{1}{\hat{\Delta}+2l_+(l_-+k_-)}\\
=& \frac{1}{\hat{\Delta}}- \frac{2(l_-+k_-)l_+}{\hat{\Delta}^2}+\frac{4(l_-+k_-)^2l_+^2}{(\hat{\Delta}+2l_+(l_-+k_-))\hat{\Delta}^2},
\end{split}
\end{align}
其中 $\hat{\Delta}=2k_+(l_-+k_-)-(\vec{l}_\perp+\vec{k}_\perp)^2$ 与 $l_+$ 无关。同样可以论证：由于 $l_+$ 被因子化出来，前两项在 $l_+$ 积分后没有贡献。然而，式~\eqref{eq:decomPlus} 的最后一项仍能从三维积分产生紫外发散。这是因为 PDF 的紫外发散来自任一圈动量按 $(l_{+}, l_{-}, \vec{l}_{\perp})\sim(1,\lambda^2,\lambda)$（当 $\lambda\to\infty$）变化的区域，于是 $l_- l_+ \sim l_\perp^2 \sim \lambda^2$。boost 不变性保证分子中的每个 $l_+$ 要么伴随分子中的一个``$-$''方向动量分量，要么伴随分母中的一个``+''方向动量分量，而这两种情形都不会压低圈动量积分的紫外发散指标。


%%%%%%%%%%%%%%%%%%%%%%%%%%%%%%%
\begin{figure}[t]
\begin{center}
\includegraphics[width=3in]{multiloopdiv}
\caption{\label{fig:multiloopdiv}
可能给夸克准 PDF 带来紫外发散贡献的三类图拓扑。}
\end{center}
\end{figure}
%%%%%%%%%%%%%%%%%%%%%%%%%%%%%%%

%%%%%%%%%%%%%%%%%%%%%%%%%%%%%%%
\subsection{夸克准 PDF 的发散图}

基于上述策略与幂次计数规则，可以从一组完备的基本构件出发找出所有紫外发散的费曼图。为了辨认夸克准 PDF 的全部紫外发散费曼图，除了图~\ref{fig:oneloop} 与图~\ref{fig:oneloopg} 中的一圈图之外，还需要图~\ref{fig:twoloop} 中的另一个图来补全基本构件集合，因为这个两圈图的表面紫外发散与在图~\ref{fig:oneloop}(b) 上按 \textbf{情形 IV} 再插入一条胶子所得的结果不一致。由于只需考虑四维积分，图~\ref{fig:twoloop} 中非零的 $\xi_z$ 保证该图是紫外有限的。

利用上述基本构件，可以生成所有可能的高阶费曼图。其中，除来自 QCD 拉氏量重正化的那些之外，只有三类拓扑可能给夸克准 PDF 带来紫外发散贡献，如图~\ref{fig:multiloopdiv} 所示。由上述讨论可知，这三类拓扑都只在所有圈动量趋于无穷的区域才可能是紫外发散的，因此紫外发散贡献在坐标空间局域。直接的推论是：要使图~\ref{fig:multiloopdiv} 中的图发散，必须有 $r_2-r_1$ 趋于 0。这一发现也解释了文献 \cite{Ji:2015jwa} 中两圈计算的结果：DR 下夸克准 PDF 的紫外发散在动量空间中正比于 $\delta(1-z)$。

除图~\ref{fig:multiloopdiv} 中的这三个紫外发散拓扑外，我们还在图~\ref{fig:multiloopfin} 中展示了一些紫外有限的拓扑。特别地，图~\ref{fig:multiloopfin} 中的最后一个图表明：在 QCD 微扰论所有阶的重正化下，夸克准 PDF 不与胶子准 PDF 混合。


%%%%%%%%%%%%%%%%%%%%%%%%%%%%%%%
\begin{figure}[h]
\begin{center}
\includegraphics[width=3in]{multiloopfin}
\caption{\label{fig:multiloopfin}
给夸克准 PDF 带来紫外有限贡献的图拓扑示例。}
\end{center}
\end{figure}
%%%%%%%%%%%%%%%%%%%%%%%%%%%%%%%

%%%%%%%%%%%%%%%%%%%%%%%%%%%%%%%
\begin{figure}[h]
\begin{center}
\includegraphics[width=3in]{sumSE}
\caption{\label{fig:sumSE}
图~\ref{fig:multiloopdiv}(a) 一般图的贡献，其中所有圈图按 1PI 图重新组合。}
\end{center}
\end{figure}
%%%%%%%%%%%%%%%%%%%%%%%%%%%%%%%
